\chapter{A complete company and a real disclosure}
\label{ch:case}

\section{What a balance must remember}

Northstar Components is a synthetic seller of one homogeneous good. Its
disclosures, transactions and contract terms are specified together, so we can
check whether reconstruction preserves a known possible history. All amounts
in this example are currency units, not thousands. The company has no foreign
currency, subsidiaries, leases, returns or other comprehensive income. These
are declared features of the example, not assumptions to impose on an imported
company. Its source package and expected results are supplied in
\code{fixtures/northstar.json}.

At the beginning of 2024, Northstar holds cash of 60, receivables of 20,
inventory of 30 and plant costing 100 with accumulated depreciation of 20.
It owes suppliers 20 and lenders 100. Share capital is 50 and retained earnings
are 20. Thus assets of 190 equal liabilities of 120 plus equity of 70.
Inventory consists of six units costing 5 each. The existing plant depreciates
by 10 a year, with no residual value. The loan remains outstanding throughout
2024; both possible contract types charge 5\% during that year.

The six months ending June and December provide two successive observation
periods. Here the chosen accrual convention assigns one half-year to each.
Cash-flow dates within a half-year are not recovered from a period aggregate;
the synthetic history supplies a feasible ordering in which inventory is
purchased before delivery and customer receipts precede the payments they fund.
More detailed liquidity questions require more detailed dates.

\begin{center}
\begin{tabular}{lrr}
\toprule Movement & First half & Second half\\\midrule
Delivered sales on credit &110&150\\
Customer collections &110&135\\
Inventory purchases on credit &70&85\\
Supplier payments &65&80\\
Cost of goods delivered &75&90\\
Payroll paid &20&22\\
Plant purchased at period end &10&12\\
Depreciation &5&6\\
Interest accrued and paid &2.5&2.5\\
Tax expense &1.5&5.9\\
Tax paid &3.5&3.9\\
Dividends paid &4&5\\\bottomrule
\end{tabular}
\end{center}

Each row has an economic meaning before it has a posting. Delivery of the
good transfers control and creates an unconditional receivable; invoicing is
assumed simultaneous in this example. Collection settles that receivable.
Purchases increase inventory and trade payables; paying the supplier settles
the payable. Delivery separately removes inventory at cost. Plant purchases
are capital expenditure, while depreciation consumes its service potential.
The addition of 10 enters service after June and has a five-year life, producing
one unit of second-half depreciation. The addition of 12 enters service after
December and has a twelve-year life.

Tax is deliberately simple: 20\% of positive profit before tax, with all these
expenses deductible and no deferred tax. The first-half payment exceeds its
tax expense by 2, leaving prepaid tax; that prepayment is used in the second
half. A genuine corporate tax package would need jurisdiction, tax basis,
losses, timing differences and payment schedules. The example's tax rule
governs its own numbers only.

\section{Recovering movements without inventing invoices}

Suppose the annual report discloses sales of 260, cost of sales of 165 and the
opening and closing working-capital balances, while the collection and purchase
totals are hidden. Under Northstar's explicit no-write-off, no-acquisition and
no-currency-effect scope, the three bridges are
\begin{align}
35&=20+260-C,\label{eq:ns-ar}\\
20&=30+P-165,\\
30&=20+P-S.
\end{align}
Here $C$ denotes collections, $P$ purchases and $S$ supplier payments, all in
currency units over the same year. Substitution gives $C=245$, $P=155$ and
$S=145$. These are identifiable aggregate movements within this model. They do
not identify which invoice was paid on which day. If acquisitions or write-offs
are admitted, their movement variables enter the same bridges and these totals
may cease to be unique.

The loan remains different. Write $F$ and $V$ for fixed and floating principal.
The source imposes $F+V=100$, $F,V\geq0$ and
$0.05F+0.05V=5$. The last equality adds no information about the split.
Both endpoints and every intermediate allocation survive. Northstar therefore
has identified working-capital totals and an unidentified financing exposure
in the same joint historical family. A solver must not turn its convenient
witness $F=0$ into the company's recovered contract.

The model's historical slots are one delivery, collection, purchase, supplier
settlement, payroll, plant addition, depreciation, interest, tax payment and
dividend per half-year. Each slot represents a classified aggregate event.
The interim filing supplies the information for allocating the annual movements
between the two half-years; the annual bridges alone identify only their totals.
The template for collection contributes $+C$ to cash and $-C$ to receivables;
its state effect reduces the unpaid receivable pool by $C$. The purchase
template adds inventory quantity $P/5$, inventory cost $P$ and a payable $P$.
A cash-balance identity by itself cannot substitute for either template.

\section{The complete historical statements}

The corrected statements follow from those movements. Temporary income and
expense accounts close to retained earnings at each half-year boundary;
annual income is the sum of the two periods' classified movements, excluding
closing entries.

\begin{center}
\begin{tabular}{lrrr}
\toprule Balance & Opening & June & December\\\midrule
Cash &60&65&74.6\\
Receivables &20&20&35\\
Inventory &30&25&20\\
Plant, net &80&85&91\\
Prepaid tax &0&2&0\\
Total assets &190&197&220.6\\\midrule
Trade payables &20&25&30\\
Debt &100&100&100\\
Share capital &50&50&50\\
Retained earnings &20&22&40.6\\
Liabilities and equity &190&197&220.6\\\bottomrule
\end{tabular}
\end{center}

Annual revenue of 260 less cost of sales 165, payroll 42, depreciation 11 and
interest 5 gives profit before tax of 37. Tax expense is 7.4 and profit is 29.6.
Other comprehensive income is zero under the declared scope, so total
comprehensive income is also 29.6. Equity moves from 70 to
$70+29.6-9=90.6$ after dividends.

The cash-flow statement supplies an independent reconciliation. Collections
245 less supplier payments 145, payroll 42, interest 5 and tax payments 7.4
produce operating cash of 45.6. Investing cash is $-22$ and financing cash
is $-9$, giving a cash increase of 14.6. The indirect route gives the same
operating cash: profit 29.6 plus depreciation 11, minus the receivables
increase 15, plus the inventory decrease 10 and the payables increase 10.
The prepaid-tax balance is zero at both annual endpoints. The two cash routes
use different information and therefore provide a useful check on classification.

\section{A correction changes what was knowable}

The original interim filing reported revenue of 120 and receivables of 30.
Ten units of revenue had been recorded before delivery. Its tax expense was
3.5, profit 14 and equity 80. The annual filing subsequently corrected the
interim revenue and receivable by 10 and recognised the excess tax payment as
prepaid tax of 2. Corrected interim profit is 6 and equity is 72, as above.
The original statement remains an accurate record of what was published, even
though its revenue recognition was wrong under the stated delivery rule.

For the temporal fixture, the original report becomes available on 1 August
2024 and is ingested on 3 August. The correction and annual report become
available on 2 March 2025 and are ingested on 5 March. An investor-information
query at 1 September 2024 00:00 UTC returns 120. For this synthetic example,
take UTC as the reporting timezone and use an instant cutoff of
3 March 2025 at 00:00 UTC, after the correction's entire publication-date
interval has ended. That query returns 110 when evaluated using the subsequently acquired archive. A query
that additionally asks what the database contained on 4 March still returns
120. Publication and ingestion answer different questions.

The annual total 260 and corrected interim total 110 share revision group
\code{v2}, so second-half revenue is $260-110=150$. Subtracting the original
120 would yield 140, but that result mixes incompatible revisions. The query
must reject the subtraction unless a reviewed compatibility relationship
supports it. A later date alone does not establish such a relationship.

\section{One economic scenario, several complete futures}

For 2025, assume demand for 30 units at price 10, unit purchase cost 5 and a
target closing inventory of eight units. Opening inventory is four units, so
the purchase decision is $30+8-4=34$ units, costing 170. Deliveries cost 150,
leaving inventory of 40. Collect 90\% of current sales and 10 of old
receivables: $270+10=280$, leaving receivables of 55. Pay old suppliers 30
and 130 of the new purchases, leaving payables of 40. Payroll is 50, capital
expenditure of 20 occurs at year-end and dividends are 8.

These decisions are inputs to a conditional forecast. They are not estimates
of demand or optimal management responses. No new borrowing is available.
Under the existing asset schedules, depreciation is $10+10/5+12/12=13$.
The new year-end asset begins depreciation in the next period. Let
$\alpha=F/100$. With the floating rate resetting to 7\%, interest is
$I=7-2\alpha$. Profit before tax is
$300-150-50-13-I=80+2\alpha$, giving tax $16+0.4\alpha$ and profit
$64+1.6\alpha$.

\begin{center}
\begin{tabular}{lrr}
\toprule 2025 statement amount & All fixed & All floating\\\midrule
Revenue &300&300\\
Interest expense &5&7\\
Tax expense &16.4&16\\
Profit and comprehensive income &65.6&64\\
Operating cash flow &48.6&47\\
Investing cash flow &$-20$&$-20$\\
Financing cash flow &$-8$&$-8$\\\midrule
Closing cash &95.2&93.6\\
Receivables / inventory / net plant &55 / 40 / 98&55 / 40 / 98\\
Assets &288.2&286.6\\
Payables / debt &40 / 100&40 / 100\\
Equity &148.2&146.6\\\bottomrule
\end{tabular}
\end{center}

The joint symbolic result is cash $93.6+1.6\alpha$, profit
$64+1.6\alpha$ and equity $146.6+1.6\alpha$, with the same
$\alpha\in[0,1]$ throughout. Selecting the smallest cash and largest equity
independently would describe no feasible company. A requested distribution
requires a measure over that same historical family; the accounting equations
do not supply one.

The later method schedules use the floating share $x$ rather than the fixed
share $\alpha$ used here: $x=1-\alpha$. Both parameterisations describe the
same joint family. Its 2025 opening cash is 74.6 and equity is 90.6; the
forecast spends 20 on period-end capital expenditure, distributes 8 and settles
its stipulated tax in the same year. These teaching assumptions are retained
when the later lease example introduces a different tax-payment schedule.

\section{Apple: source labels do not identify exposure}
\label{sec:apple}

The real reference package is Apple's 2024 Form 10-K and first-quarter 2025
Form 10-Q, for fiscal dates 28 September and 28 December 2024 respectively
\citep{apple2024,apple2025q1}. The annual balance sheet also supplies
30 September 2023. The annual filing is archived as the company's PDF from
AnnualReports; the interim filing is archived from SEC EDGAR with its inline
XBRL. Apple's Note 1 describes a 52/53-week year ending on the last Saturday
of September. Calendar-year assumptions would therefore misstate the periods.

The package inventory preserves every annual PDF page and every interim
top-level section, all thirteen annual notes and all ten interim notes.
It records internal cross-references and externally incorporated material,
including the proxy statement, as distinct obligations. Statement and debt
facts below are structured; policies, other notes and exhibits retain their
source locations and specific modelling obligations. This is complete
inventory coverage of the acquired files, with partial semantic formalisation.
It does not claim to reconstruct every Apple contract or to have acquired
every incorporated external document.

Apple's annual balance sheet reports assets 364,980, liabilities 308,030 and
equity 56,950, in millions of US dollars. Note 9 reports term-debt principal
97,341, unamortised discount/premium and issuance costs of $-321$, and hedge
accounting adjustments of $-358$. Thus its carrying amount is
$97{,}341-321-358=96{,}662$, split into current 10,912 and non-current
85,750. Fair value is separately disclosed as 88.4 billion. None of these
three measurements may replace another merely because a table calls it debt.

Note 9 identifies the notes as fixed-rate, then states that interest-rate
swaps effectively convert a portion to floating rates. Note 4 reports
interest-rate contract notional of 14,575 at the annual date. A notional is
neither a liability carrying amount nor sufficient information for an exact
interest sensitivity: reset timing, remaining periods, curve conventions and
hedge coverage still matter. The cross-reference connects the derivative to
the borrowing without erasing either contract.

At the interim date, total term-debt carrying amount is
$10{,}848+83{,}956=94{,}804$. Cash is 30,299, total assets 344,085,
liabilities 277,327 and equity 66,758. The interim debt note's rounded
``94.8 billion'' and the statement's 94,804 million can coexist; their
precision differs. Treating them as conflicting exact numbers would create
an artificial infeasibility.

To demonstrate a source-to-forecast calculation, freeze the interim balance
sheet and apply a \emph{marginal} financing stress for a following 13-week
period: a 200-basis-point rate increase, no principal change, no offsetting
hedge beyond the explicitly selected exposure, immediate cash settlement and
no tax effect. An assumed effective floating exposure $E=14{,}575$ million
produces $0.02\times(13/52)E=72.875$ million additional interest. This exposure
uses the annual notional as a conditional input; it is not asserted to be the
interim net sensitivity. Under that condition, debit interest expense and
credit cash by 72.875. The resulting marginal statements have cash 30,226.125,
assets 344,012.125, liabilities 277,327, equity 66,685.125, profit and comprehensive
income of $-72.875$, and operating cash flow of $-72.875$ relative to the frozen
baseline. The balance-sheet difference equals the equity difference exactly.

If exposure is instead only restricted to $0\leq E\leq14{,}575$, the
conditional added interest lies in $[0,72.875]$ before monetary rounding. That
restriction is a stress assumption and is not claimed to bound every real
Apple exposure. This real case verifies evidence, context, mappings and a
coherent marginal propagation. Northstar supplies the complete-company
forecast; a complete Apple forecast remains dependent on its unmodelled
operations and undisclosed terms.
